\textbf{Phần 1. Trắc nghiệm lựa chọn}
\setcounterpageref{ex}{0}

\begin{ex} %[5D6N35]
Hỗn số $5\dfrac{2}{5}$ được viết dưới dạng phân số là
\choice
{$\dfrac{10}{5}$}
{\True $\dfrac{27}{5}$}
{$\dfrac{12}{5}$}
{$\dfrac{52}{5}$}
\loigiai{
Ta có: $5\dfrac{2}{5} = \dfrac{5 \times 5 + 2}{5} = \dfrac{27}{5}$.\\
\textbf{Đáp án của câu hỏi là B}
}
\end{ex}

\begin{ex} %[5D6N35]
Phân số $\dfrac{96}{50}$ được viết dưới dạng số thập phân là
\choice
{\True $1{,}92$}
{$0{,}96$}
{$9{,}6$}
{$19{,}2$}
\loigiai{
Ta có: $\dfrac{96}{50} = \dfrac{192}{100} = 1{,}92$.\\
\textbf{Đáp án của câu hỏi là A}
}
\end{ex}

\begin{ex} %[5D6H35]
Kết quả của biểu thức tính nhanh $\dfrac{2}{7} \times \dfrac{4}{9} + \dfrac{2}{7} \times \dfrac{5}{9}$ là
\choice
{$\dfrac{4}{7}$}
{$\dfrac{18}{63}$}
{$\dfrac{9}{16}$}
{\True $\dfrac{2}{7}$}
\loigiai{
Ta có: $\dfrac{2}{7} \times \left(\dfrac{4}{9} + \dfrac{5}{9}\right) = \dfrac{2}{7} \times 1 = \dfrac{2}{7}$.\\
\textbf{Đáp án của câu hỏi là D}
}
\end{ex}

\begin{ex} %[5D6H35]
Số gồm: \textit{Ba mươi lăm đơn vị, hai phần mười, không phần trăm, sáu phần nghìn} được viết là
\choice
{$35{,}26$}
{$35{,}026$}
{\True $35{,}206$}
{$35{,}260$}
\loigiai{
Số thập phân đó là $35{,}206$.\\
\textbf{Đáp án của câu hỏi là C}
}
\end{ex}

\begin{ex} %[5D6H35]
Đổi $1\,052\text{ ha}$ ra ki-lô-mét vuông ta được số thập phân là
\choice
{$0{,}1052\text{ km}^2$}
{$1{,}052\text{ km}^2$}
{$105{,}2\text{ km}^2$}
{\True $10{,}52\text{ km}^2$}
\loigiai{
Vì $1\text{ km}^2 = 100\text{ ha}$ nên $1\,052\text{ ha} = 1\,052 : 100 = 10{,}52\text{ km}^2$.\\
\textbf{Đáp án của câu hỏi là D}
}
\end{ex}

\begin{ex} %[5D6H35]
Tính giá trị của biểu thức $12{,}5 \times 3{,}6 + 12{,}5 \times 2{,}4$ ta được
\choice
{$60$}
{$750$}
{\True $75$}
{$150$}
\loigiai{
Ta có: $12{,}5 \times (3{,}6 + 2{,}4) = 12{,}5 \times 6 = 75$.\\
\textbf{Đáp án của câu hỏi là C}
}
\end{ex}

\begin{ex} %[5D6V35]
Trong đợt thu gom giấy vụn, lớp 5A gom được $45\text{ kg}$, lớp 5B gom được bằng $\dfrac{2}{3}$ lớp 5A, lớp 5C gom ít hơn lớp 5B là $15\text{ kg}$. Biết $1\text{ kg}$ giấy vụn làm được $15$ cuốn vở tái chế. Số cuốn vở tái chế cả ba lớp làm được là
\choice
{\True $1\,350$ cuốn}
{$90$ cuốn}
{$675$ cuốn}
{$1\,500$ cuốn}
\loigiai{
Lớp 5B gom được: $45 \times \dfrac{2}{3} = 30\text{ (kg)}$.\\
Lớp 5C gom được: $30 - 15 = 15\text{ (kg)}$.\\
Cả ba lớp gom được: $45 + 30 + 15 = 90\text{ (kg)}$.\\
Số cuốn vở tái chế làm được là: $90 \times 15 = 1\,350$ (cuốn).\\
\textbf{Đáp án của câu hỏi là A}
}
\end{ex}

\begin{ex} %[5D6V35]
Có $6\text{ kg}$ đường chia đều vào $12$ gói. Nếu $1\text{ kg}$ đường có giá $17\,000$ đồng thì người mua $10$ gói đường như vậy phải trả số tiền là
\choice
{$34\,000\text{ đồng}$}
{\True $85\,000\text{ đồng}$}
{$68\,000\text{ đồng}$}
{$170\,000\text{ đồng}$}
\loigiai{
Mỗi gói đường nặng: $6 : 12 = 0{,}5\text{ (kg)}$.\\
$10$ gói đường nặng: $0{,}5 \times 10 = 5\text{ (kg)}$.\\
Số tiền phải trả là: $5 \times 17\,000 = 85\,000$ (đồng).\\
\textbf{Đáp án của câu hỏi là B}
}
\end{ex}

\begin{ex} %[5H6V35]
Cho hình tam giác vuông $ABC$ có hai cạnh góc vuông là $6\text{ cm}$ và $8\text{ cm}$, cạnh huyền là $10\text{ cm}$. Tam giác nằm trong hình tròn đường kính $10\text{ cm}$. Diện tích của hình tròn đó là
\choice
{$31{,}4\text{ cm}^2$}
{$314\text{ cm}^2$}
{$24\text{ cm}^2$}
{\True $78{,}5\text{ cm}^2$}
\loigiai{
Bán kính hình tròn là: $10 : 2 = 5\text{ (cm)}$.\\
Diện tích hình tròn là: $3{,}14 \times 5 \times 5 = 78{,}5\text{ (cm}^2\text{)}$.\\
\textbf{Đáp án của câu hỏi là D}
}
\end{ex}

\begin{ex} %[5H6V35]
Cho hình thang vuông $ABCD$ có đáy lớn hơn đáy bé $40\text{ cm}$, chiều cao bằng trung bình cộng hai đáy và bằng $80\text{ cm}$. Đoạn thẳng $BE$ vuông góc với $CD$ chia hình thang thành hình chữ nhật $ABED$ và tam giác vuông $BCE$. Diện tích tam giác $BCE$ là
\choice
{$3\,200\text{ cm}^2$}
{\True $1\,600\text{ cm}^2$}
{$800\text{ cm}^2$}
{$4\,800\text{ cm}^2$}
\loigiai{
Đáy $EC$ của tam giác $BCE$ chính bằng hiệu độ dài hai đáy: $EC = 40\text{ cm}$.\\
Chiều cao $BE = 80\text{ cm}$.\\
Diện tích tam giác $BCE$ là: $\dfrac{40 \times 80}{2} = 1\,600\text{ (cm}^2\text{)}$.\\
\textbf{Đáp án của câu hỏi là B}
}
\end{ex}

\begin{ex} %[5D6V35]
Giá trị của biểu thức $(64{,}2 - 36{,}6) : 1{,}2 + 13{,}15$ là
\choice
{$23$}
{$35{,}15$}
{\True $36{,}15$}
{$46{,}15$}
\loigiai{
$(64{,}2 - 36{,}6) : 1{,}2 + 13{,}15 = 27{,}6 : 1{,}2 + 13{,}15 = 23 + 13{,}15 = 36{,}15$.\\
\textbf{Đáp án của câu hỏi là C}
}
\end{ex}

\begin{ex} %[5H6V35]
Một hình thang vuông $ABCD$ có $AB = 16\text{ cm}$, $CD = 28\text{ cm}$ và chiều cao $AD = \dfrac{5}{7}CD$. Bên trong hình thang có một nửa hình tròn đường kính $AD$. Diện tích phần hình thang nằm ngoài nửa hình tròn là
\choice
{\True $283\text{ cm}^2$}
{$440\text{ cm}^2$}
{$157\text{ cm}^2$}
{$314\text{ cm}^2$}
\loigiai{
Chiều cao $AD = 28 \times \dfrac{5}{7} = 20\text{ (cm)}$.\\
Diện tích hình thang: $\dfrac{(16 + 28) \times 20}{2} = 440\text{ (cm}^2\text{)}$.\\
Bán kính nửa hình tròn là $20 : 2 = 10\text{ (cm)}$.\\
Diện tích nửa hình tròn là: $3{,}14 \times 10 \times 10 : 2 = 157\text{ (cm}^2\text{)}$.\\
Diện tích phần tô màu còn lại: $440 - 157 = 283\text{ (cm}^2\text{)}$.\\
\textbf{Đáp án của câu hỏi là A}
}
\end{ex}

\textbf{Phần 2. Trắc nghiệm đúng sai}
\setcounterpageref{ex}{0}

\begin{ex} %[5D6H35]
Xét tính đúng sai của các khẳng định sau:
\choiceTF
{\True Phân số $\dfrac{327}{300}$ viết thành số thập phân là $1{,}09$.}
{Số $5{,}099$ lớn hơn số $5{,}1$.}
{\True Chữ số $8$ trong số $30{,}812$ ở hàng phần mười.}
{\True Phép tính $75{,}52 : 16 = 4{,}72$.}
\loigiai{
\begin{itemize}
    \item $\dfrac{327}{300} = \dfrac{109}{100} = 1{,}09$.
    \item $5{,}099 < 5{,}100 = 5{,}1$ (phần mười có $0 < 1$). Khẳng định lớn hơn là sai.
    \item Chữ số 8 đứng ngay sau dấu phẩy nên ở hàng phần mười.
    \item $75{,}52 : 16 = 4{,}72$.
\end{itemize}
\begin{itemize}
    \item a) \textbf{Đúng}.
    \item b) \textbf{Sai}.
    \item c) \textbf{Đúng}.
    \item d) \textbf{Đúng}.
\end{itemize}
}
\end{ex}

\begin{ex} %[5H6V35]
Cho hình thang vuông $ABCD$ có đáy bé $AB = 16\text{ cm}$, đáy lớn $CD = 28\text{ cm}$, chiều cao $AD = 20\text{ cm}$ và một nửa hình tròn đường kính $AD$ nằm bên trong hình thang. Xét tính đúng sai của các khẳng định sau:
\choiceTF
{\True Diện tích hình thang vuông $ABCD$ là $440\text{ cm}^2$.}
{\True Bán kính của nửa hình tròn là $10\text{ cm}$.}
{Diện tích của nửa hình tròn là $314\text{ cm}^2$.}
{\True Diện tích phần hình thang nằm ngoài nửa hình tròn là $283\text{ cm}^2$.}
\loigiai{
\begin{itemize}
    \item $S_{thang} = \dfrac{(16 + 28) \times 20}{2} = 440\text{ cm}^2$.
    \item Bán kính $r = 20 : 2 = 10\text{ cm}$.
    \item Diện tích cả hình tròn là $314\text{ cm}^2 \Rightarrow$ Nửa hình tròn là $157\text{ cm}^2 \neq 314\text{ cm}^2$.
    \item Diện tích phần tô màu: $440 - 157 = 283\text{ cm}^2$.
\end{itemize}
\begin{itemize}
    \item a) \textbf{Đúng}.
    \item b) \textbf{Đúng}.
    \item c) \textbf{Sai}: Diện tích nửa hình tròn là $157\text{ cm}^2$.
    \item d) \textbf{Đúng}.
\end{itemize}
}
\end{ex}

\textbf{Phần 3. Trả lời ngắn}
\setcounterpageref{ex}{0}

\begin{ex} %[5D6H35]
Viết phân số $\dfrac{204}{125}$ dưới dạng số thập phân.
\shortans{1,632}
\loigiai{
\textbf{Đáp số:} $1{,}632$\\
$\dfrac{204}{125} = \dfrac{1\,632}{1\,000} = 1{,}632$.
}
\end{ex}

\begin{ex} %[5D6H35]
Tính giá trị biểu thức: $12{,}5 \times 3{,}6 + 12{,}5 \times 2{,}4$.
\shortans{75}
\loigiai{
\textbf{Đáp số:} $75$\\
$12{,}5 \times (3{,}6 + 2{,}4) = 12{,}5 \times 6 = 75$.
}
\end{ex}

\begin{ex} %[5D6H35]
Tính kết quả phép tính: $(64{,}2 - 36{,}6) : 1{,}2 + 13{,}15$. (Nhập kết quả dưới dạng số thập phân).
\shortans{36,15}
\loigiai{
\textbf{Đáp số:} $36{,}15$\\
$27{,}6 : 1,2 + 13,15 = 23 + 13,15 = 36,15$.
}
\end{ex}

\begin{ex} %[5D6V35]
Tổng số giấy vụn $3$ lớp $5A, 5B, 5C$ thu gom được là $90\text{ kg}$. Biết $1\text{ kg}$ giấy vụn làm được $15$ cuốn vở tái chế. Hỏi làm được tất cả bao nhiêu cuốn vở tái chế?
\shortans{1350}
\loigiai{
\textbf{Đáp số:} $1350$\\
$90 \times 15 = 1\,350$ (cuốn vở).
}
\end{ex}

\begin{ex} %[5H6V35]
Cho hình thang vuông $ABCD$ có đáy lớn $CD = 100\text{ cm}$, đáy bé $AB = 60\text{ cm}$, chiều cao $80\text{ cm}$. Kẻ $BE \perp CD$. Tính diện tích tam giác $BCE$ theo xăng-ti-mét vuông.
\shortans{1600}
\loigiai{
\textbf{Đáp số:} $1600$\\
$\dfrac{(100 - 60) \times 80}{2} = 1\,600\text{ (cm}^2\text{)}$.
}
\end{ex}

\begin{ex} %[5H6V35]
Hình thang vuông có hai đáy $16\text{ cm}, 28\text{ cm}$ và chiều cao $20\text{ cm}$. Nửa hình tròn bên trong có đường kính $20\text{ cm}$. Tính diện tích phần hình thang nằm ngoài nửa hình tròn theo xăng-ti-mét vuông.
\shortans{283}
\loigiai{
\textbf{Đáp số:} $283$\\
$\dfrac{(16 + 28) \times 20}{2} - \dfrac{3{,}14 \times 10 \times 10}{2} = 440 - 157 = 283\text{ (cm}^2\text{)}$.
}
\end{ex}
